\begin{ex}
Số bạn yêu thích môn thể thao của lớp 3A: Bóng đá: 18 bạn, Cầu lông: 12 bạn, Bơi lội: 10 bạn. Môn nào được nhiều bạn thích nhất? (Bóng đá ghi 1, Cầu lông ghi 2, Bơi lội ghi 3).
\shortans{1}
\loigiai{
Bóng đá có 18 bạn.
}
\end{ex}

\begin{ex}
Tổng số học sinh tham gia bình chọn môn thể thao trên là bao nhiêu bạn?
\shortans{40}
\loigiai{
$18 + 12 + 10 = 40$ bạn.
}
\end{ex}

\begin{ex}
Số bạn thích bóng đá nhiều hơn số bạn thích bơi lội bao nhiêu bạn?
\shortans{8}
\loigiai{
$18 - 10 = 8$ bạn.
}
\end{ex}

\begin{ex}
Bảng số cây trồng của 3 lớp: 3A: 45 cây, 3B: 50 cây, 3C: 40 cây. Cả 3 lớp trồng được bao nhiêu cây?
\shortans{135}
\loigiai{
$45 + 50 + 40 = 135$ cây.
}
\end{ex}

\begin{ex}
Lớp trồng được nhiều cây nhất trồng bao nhiêu cây?
\shortans{50}
\loigiai{
Lớp 3B trồng 50 cây.
}
\end{ex}

\begin{ex}
Gieo một con xúc xắc 6 mặt 1 lần. Khả năng xuất hiện mặt 8 chấm có thể xảy ra không? (Có ghi 1, Không ghi 0).
\shortans{0}
\loigiai{
Không thể có mặt 8 chấm.
}
\end{ex}

\begin{ex}
Gieo một đồng xu 1 lần. Có mấy kết quả có thể xảy ra? (Sấp, Ngửa).
\shortans{2}
\loigiai{
Có 2 khả năng.
}
\end{ex}

\begin{ex}
Trong hộp có 6 bi xanh, không có bi nào khác. Lấy ngẫu nhiên 1 viên bi thì chắc chắn lấy được bi màu gì? (Bi xanh ghi 1, Bi đỏ ghi 2).
\shortans{1}
\loigiai{
Chắc chắn là bi xanh.
}
\end{ex}

\begin{ex}
Trong hộp có 5 viên bi đỏ và 5 viên bi vàng. Lấy ngẫu nhiên 1 viên bi, có thể lấy được bi xanh không? (Có ghi 1, Không thể ghi 0).
\shortans{0}
\loigiai{
Không thể lấy được bi xanh.
}
\end{ex}

\begin{ex}
Tung xúc xắc 1 lần, sự kiện xuất hiện mặt có số chấm bé hơn 7 là sự kiện gì? (Chắc chắn ghi 1, Có thể ghi 2, Không thể ghi 3).
\shortans{1}
\loigiai{
Chắc chắn (tất cả các mặt từ 1 đến 6 đều < 7).
}
\end{ex}

\begin{ex}
Tung xúc xắc 1 lần, sự kiện xuất hiện mặt 3 chấm là sự kiện gì? (Chắc chắn ghi 1, Có thể ghi 2, Không thể ghi 3).
\shortans{2}
\loigiai{
Có thể xảy ra.
}
\end{ex}

\begin{ex}
Tung xúc xắc 1 lần, sự kiện xuất hiện mặt 0 chấm là sự kiện gì? (Chắc chắn ghi 1, Có thể ghi 2, Không thể ghi 3).
\shortans{3}
\loigiai{
Không thể xảy ra.
}
\end{ex}

\begin{ex}
Trong hộp có 7 viên bi đỏ và 3 viên bi xanh. Muốn chắc chắn lấy được ít nhất 1 viên bi xanh thì cần lấy ra ít nhất bao nhiêu viên bi?
\shortans{8}
\loigiai{
Trường hợp xấu nhất lấy hết 7 bi đỏ, lấy thêm 1 viên: $7 + 1 = 8$ viên.
}
\end{ex}

\begin{ex}
Trong túi có 6 quả cầu trắng và 4 quả cầu đen. Cần lấy ít nhất mấy quả cầu để chắc chắn có quả cầu trắng?
\shortans{5}
\loigiai{
$4 + 1 = 5$ quả.
}
\end{ex}

\begin{ex}
Số sách đọc trong tháng của 4 bạn: An: 4 cuốn, Bình: 6 cuốn, Chi: 5 cuốn, Dung: 7 cuốn. Ai đọc được nhiều sách nhất? (An: 1, Bình: 2, Chi: 3, Dung: 4).
\shortans{4}
\loigiai{
Dung đọc 7 cuốn (nhiều nhất).
}
\end{ex}

\begin{ex}
Tổng số sách cả 4 bạn đọc được là bao nhiêu cuốn?
\shortans{22}
\loigiai{
$4 + 6 + 5 + 7 = 22$ cuốn.
}
\end{ex}

\begin{ex}
Bình đọc nhiều hơn An bao nhiêu cuốn sách?
\shortans{2}
\loigiai{
$6 - 4 = 2$ cuốn.
}
\end{ex}

\begin{ex}
Tung đồng xu 30 lần, có 16 lần xuất hiện mặt Sấp. Có bao nhiêu lần xuất hiện mặt Ngửa?
\shortans{14}
\loigiai{
$30 - 16 = 14$ lần.
}
\end{ex}

\begin{ex}
Quay vòng quay có 6 ô số từ 1 đến 6. Có bao nhiêu khả năng mũi tên dừng lại ở ô số chẵn?
\shortans{3}
\loigiai{
Có 3 ô số chẵn: 2, 4, 6.
}
\end{ex}

\begin{ex}
Quay vòng quay có 6 ô số từ 1 đến 6. Có bao nhiêu khả năng dừng lại ở ô số lẻ?
\shortans{3}
\loigiai{
Có 3 ô số lẻ: 1, 3, 5.
}
\end{ex}

\begin{ex}
Hai bạn cùng chơi tung xúc xắc. Tổng số chấm lớn nhất có thể nhận được khi tung 2 con xúc xắc là bao nhiêu?
\shortans{12}
\loigiai{
$6 + 6 = 12$ chấm.
}
\end{ex}

\begin{ex}
Tổng số chấm nhỏ nhất có thể nhận được khi tung 2 con xúc xắc là bao nhiêu?
\shortans{2}
\loigiai{
$1 + 1 = 2$ chấm.
}
\end{ex}

\begin{ex}
Kiểm đếm số học sinh đạt điểm 10 môn Tiếng Việt của khối 3: Lớp 3A: 8 bạn, Lớp 3B: 12 bạn, Lớp 3C: 10 bạn. Cả khối có bao nhiêu bạn đạt điểm 10?
\shortans{30}
\loigiai{
$8 + 12 + 10 = 30$ bạn.
}
\end{ex}

\begin{ex}
Lớp 3B nhiều hơn lớp 3A bao nhiêu bạn đạt điểm 10?
\shortans{4}
\loigiai{
$12 - 8 = 4$ bạn.
}
\end{ex}

\begin{ex}
Trong một tuần lễ có bao nhiêu ngày?
\shortans{7}
\loigiai{
Có 7 ngày.
}
\end{ex}
